\section{Core Infrastructure and Technology Rationale}\label{sec:core}
% Rubric criterion 4 (20 %): justification of OSPF, BGP, EVPN-VXLAN, STP, HSRP; underlay/overlay interaction.

Each subsection follows the same pattern: the requirement, the choice and the alternatives we rejected,
the configuration, and the behaviour on failure. Listings are excerpts of the generated startup-configs;
the complete files are in \texttt{Configurations/}.

\subsection{Layer 2: VLANs, trunks, EtherChannel and spanning tree}\label{sec:l2}
\textbf{Root placement follows the gateway.} At HQ both core switches carry every VLAN. For each VLAN we
pin the spanning-tree root to the core switch that is also the active HSRP gateway of that VLAN
(\cref{lst:stp}): \dev{HCMC-CORE-SW01} is root for VLANs 10, 30, 50, 90, 99 and 100 (priority 4096) and
secondary for 20, 40 and 60 (8192); \dev{HCMC-CORE-SW02} is the mirror image. A frame from an access switch
therefore travels up the forwarding uplink to the switch that routes it and never crosses Po1 to reach
its gateway. Splitting the VLAN groups also uses both uplinks of every access switch -- each uplink
forwards for one group and blocks for the other -- instead of leaving one idle.

\cfgsnippet[lst:stp]{l2_stp_core.cfg}{Deterministic RSTP root placement on \dev{HCMC-CORE-SW01}}

We kept Rapid PVST+ \cite{ieee8021d}: with fewer than a dozen VLANs per site the per-VLAN instances cost
nothing and converge in about a second. MST becomes the better choice when the VLAN count grows into the
hundreds, because it maps many VLANs to a few instances. Loop guard is enabled globally so that a port that
stops receiving BPDUs (a unidirectional link) blocks instead of forwarding.

\textbf{EtherChannel between the cores.} The two cores are joined by a two-member port-channel negotiated
with LACP \cite{ieee8021ax} (\cref{lst:po1}). Losing one member costs bandwidth but triggers no
spanning-tree or routing change. LACP was chosen over a static bundle (\texttt{mode on}) because it detects a
mis-cabled or one-way member and removes it, where a static bundle would keep sending traffic into it. The
trunk uses a dedicated native VLAN (998) that carries no hosts, disables DTP (\texttt{nonegotiate}) and
allows only the VLANs that exist, which removes the classic VLAN-hopping vectors \cite{ieee8021q}.

\cfgsnippet[lst:po1]{l2_portchannel.cfg}{LACP port-channel between the core switches}

\textbf{Root bridge enforcement.} A low priority makes a switch the preferred root, but does not prevent a
switch plugged in with priority~0 from taking over. The core ports towards the access layer therefore run
\emph{root guard} (\cref{lst:rootguard}): if a superior BPDU arrives, the port is put into the
root-inconsistent state instead of re-electing the root. On the access side, every user port has PortFast
and BPDU guard, so a port where someone connects a switch is shut down (\cref{sec:security}).

\cfgsnippet[lst:rootguard]{l2_trunk_rootguard.cfg}{Core-to-access trunk with root guard}

\subsection{First-hop redundancy and centralised IP services}\label{sec:fhrp}
\textbf{HSRPv2.} Every host uses the virtual address \texttt{.1} of its subnet as default gateway; two
physical gateways share it. We chose HSRP version~2 \cite{rfc2281} because the whole lab is Cisco and it
offers everything we need on IOS: IPv6 groups, millisecond timers, MD5 authentication and object tracking.
VRRPv3 \cite{rfc5798} provides the same function and would replace HSRP in a multi-vendor network; GLBP
load-balances per host but adds a forwarder election we do not need, since we already balance per VLAN.
Every IPv4 group has an IPv6 twin with group number \(+1000\) and identical settings -- an HSRP group
serves one address family -- and \cref{tab:hsrp} in the appendix lists the owner of each of the
\FactNumHsrpGroups{} IPv4 groups.

\cfgsnippet[lst:hsrp-active]{hsrp_core_svi_peer.cfg}{HSRPv2 gateway of VLAN 20, active side (\dev{HCMC-CORE-SW02})}

\Cref{lst:hsrp-active} shows the active side of VLAN~20 at HQ. Four details carry the design:

\begin{itemize}
  \item \textbf{Timers} of 500\,ms hello and 1.5\,s hold detect a dead gateway in at most 1.5\,s instead of
        the default 10\,s.
  \item \textbf{Preempt with a delay only on the designed owner.} After a reboot, the owner waits 60\,s
        before reclaiming the gateway, so it does not attract traffic before its routing has converged.
        Its partner preempts without delay, so it takes over at once when the owner's priority drops.
  \item \textbf{Tracking.} Track object~1 follows the route that makes the gateway useful: the default route
        on the HQ cores, the route to the HQ aggregate \texttt{10.1.0.0/16} on the branch routers. If the
        route disappears, the priority falls by 30 -- more than the 10-point gap -- and the partner becomes
        active. \texttt{track timer ip route 1} polls the route every second instead of every 15\,s.
  \item \textbf{Load sharing.} Active roles alternate per VLAN between the two gateways of a site, so both
        uplinks and both providers carry traffic in normal operation; a backup path that is used every day
        is known to work.
\end{itemize}

Tracking matters most at the branches (\cref{lst:hsrp-branch}). When a branch router loses its provider, its
LAN interfaces stay up, so plain HSRP would keep it active and silently drop every packet for HQ. Because it
tracks the BGP route to HQ, it hands its VLANs to the other router, whose tunnel runs over the other
provider.

\cfgsnippet[lst:hsrp-branch]{hsrp_branch_subif.cfg}{Branch gateway with route tracking (\dev{HANOI-EDGE-RTR01}, standby for VLAN 20)}

\textbf{Centralised DHCP, DNS and NTP.} As the brief requires, DHCP is centralised at HQ: the server
\dev{HCMC-DC-SRV01} (\texttt{10.1.100.10}) serves all three sites, and every gateway relays with \texttt{ip helper-address} and, for IPv6,
\texttt{ipv6 dhcp relay destination}. The relay writes its own interface address into the
\texttt{giaddr} field of the request, which is how the server selects the pool of the right subnet at the
right site \cite{rfc2131}; the reply returns to that address -- through the tunnel for a branch. Each pool
hands out the HSRP virtual address as default router (\cref{lst:dhcp}). IPv6 hosts form their address with
SLAAC and obtain DNS through stateless DHCPv6 (\texttt{other-config-flag}). The same node is authoritative
for \texttt{tvi.local} and is the NTP master for all devices, so log timestamps from all three sites can be
correlated.

\cfgsnippet[lst:dhcp]{dhcp_pool.cfg}{Central DHCP pool for a Hanoi VLAN (\dev{HCMC-DC-SRV01})}

\subsection{Underlay routing: OSPFv2 and OSPFv3}\label{sec:underlay}
\textbf{One OSPF domain per site, area 0 only.} The brief prescribes OSPF inside a site. Each site runs its
own OSPF process (OSPFv2 \cite{rfc2328} for IPv4, OSPFv3 \cite{rfc5340} for IPv6) in a single backbone area.
With at most four routers per site, more areas would add ABRs and summarisation work without reducing
anything. The sites do \emph{not} share an IGP: an interface flap in Hanoi does not cause an SPF calculation
in HCMC, and the choice of WAN path is expressed with BGP attributes, not with OSPF metrics.

\textbf{Adjacencies only where routers meet.} \texttt{passive-interface default} keeps OSPF silent on all
user VLANs, so a PC running a routing daemon cannot inject routes (\cref{lst:ospf-core}). On a core switch,
only the transit VLAN to the other core and the two uplinks to the edge routers form adjacencies. All
adjacency links use the point-to-point network type (no DR election) and a hello/dead interval of 1\,s/4\,s:
in a virtual lab -- and across any intermediate Layer-2 device -- a dead neighbour does not take the local
interface down, so the default 40\,s dead interval would be the outage. The reference bandwidth of
10\,Gbit/s keeps 1, 10 and 40\,Gbit/s links distinguishable in the metric.

\cfgsnippet[lst:ospf-core]{ospf_core.cfg}{OSPF and OSPFv3 on a core switch}

\textbf{A conditional default route.} The edge routers inject the default route into OSPF as an external
type-1 route with seed metric 10 on \dev{HCMC-EDGE-RTR01} and 20 on \dev{HCMC-EDGE-RTR02}
(\cref{lst:ospf-default}). Type~1 adds the internal path cost, so the choice is deterministic. The IPv4
statement has no \texttt{always} keyword: an edge router advertises the default only while it has one itself,
learned by eBGP from its provider. When ISP-A fails, edge router~1 withdraws its default and the cores
switch to edge router~2 without any other mechanism. (In IPv6 the providers offer no route in the lab, so
the IPv6 default only leads to the edge routers, which hold the inter-site routes; it is originated
unconditionally.)

\cfgsnippet[lst:ospf-default]{ospf_edge_default.cfg}{Conditional default origination on the HQ edge}

\textbf{The edge-to-edge link.} The core switches run no BGP and know only ``default $\rightarrow$ edge''. If
edge router~1 has to send traffic for a branch to edge router~2 (for example because its own tunnel to that
branch is down), a path through a core would fail: the core does not know the branch prefix and returns the
packet to edge router~1 on its default route -- a loop. The direct routed link between the two edge routers
(\cref{lst:edge-link}) keeps every BGP next hop one hop away; we found this case with the routing
emulation (\cref{sec:challenges}).

\subsection{Overlay and WAN edge: eBGP, IPsec and iBGP route reflection}\label{sec:overlay}
\textbf{How underlay and overlay work together.} OSPF carries everything a site needs internally: the
loopbacks of its routers, its links, its user subnets and a default route. BGP carries what lies outside:
three site aggregates between the sites, and a default route from each provider. The two meet only at the
edge routers. A packet from an HQ PC to Hanoi is routed by its gateway along the OSPF default to an HQ edge
router; there BGP knows \texttt{10.2.0.0/16} with the tunnel as next hop; the tunnel itself is carried by
the provider network to the branch, where the branch router delivers the packet into its own OSPF domain.
\Cref{tab:bgp} lists the sessions.

% GENERATED by Tools/gen_latex.py - do not edit; change Tools/tvi_model.py and regenerate.
\begin{table}[htbp]
\centering\scriptsize
\caption{BGP sessions (IPv4; every iBGP session has an IPv6 twin)}\label{tab:bgp}
\begin{tabularx}{\textwidth}{l l l l l l X}
\toprule
\textbf{Router} & \textbf{AS} & \textbf{Peer} & \textbf{Peer AS} & \textbf{Type} & \textbf{Transport} & \textbf{Role} \\
\midrule
\texttt{DANANG-EDGE-RTR01} & AS 65000 & \texttt{ISP-A-RTR} & AS 64501 & eBGP & direct &  \\
\texttt{DANANG-EDGE-RTR01} & AS 65000 & \texttt{HCMC-EDGE-RTR01} & AS 65000 & iBGP & Tunnel3 &  \\
\texttt{DANANG-EDGE-RTR02} & AS 65000 & \texttt{ISP-B-RTR} & AS 64502 & eBGP & direct &  \\
\texttt{DANANG-EDGE-RTR02} & AS 65000 & \texttt{HCMC-EDGE-RTR02} & AS 65000 & iBGP & Tunnel4 &  \\
\texttt{HCMC-EDGE-RTR01} & AS 65000 & \texttt{ISP-A-RTR} & AS 64501 & eBGP & direct &  \\
\texttt{HCMC-EDGE-RTR01} & AS 65000 & \texttt{HCMC-EDGE-RTR02} & AS 65000 & iBGP & Loopback0 &  \\
\texttt{HCMC-EDGE-RTR01} & AS 65000 & \texttt{HANOI-EDGE-RTR01} & AS 65000 & iBGP & Tunnel1 & RR client \\
\texttt{HCMC-EDGE-RTR02} & AS 65000 & \texttt{ISP-B-RTR} & AS 64502 & eBGP & direct &  \\
\texttt{HCMC-EDGE-RTR02} & AS 65000 & \texttt{HANOI-EDGE-RTR02} & AS 65000 & iBGP & Tunnel2 & RR client \\
\texttt{HANOI-EDGE-RTR01} & AS 65000 & \texttt{ISP-A-RTR} & AS 64501 & eBGP & direct &  \\
\texttt{HANOI-EDGE-RTR02} & AS 65000 & \texttt{ISP-B-RTR} & AS 64502 & eBGP & direct &  \\
\texttt{INTERNET-RTR} & AS 64510 & \texttt{ISP-A-RTR} & AS 64501 & eBGP & direct &  \\
\texttt{INTERNET-RTR} & AS 64510 & \texttt{ISP-B-RTR} & AS 64502 & eBGP & direct &  \\
\bottomrule
\end{tabularx}
\end{table}


\textbf{eBGP to the providers.} Each edge router peers with its provider (MD5-authenticated, 3\,s/9\,s
timers). Inbound, a prefix list accepts only \texttt{0.0.0.0/0}; outbound, \texttt{PL-NONE-OUT} announces
nothing, because TVI's addresses are private and leave the company only through NAT. TVI can therefore
never become a transit path between the providers \cite{rfc4271}.

\textbf{GRE over IPsec.} The tunnels are GRE interfaces protected by an IPsec profile
(\cref{lst:ipsec}): IKEv1 with AES-256, SHA-256 and Diffie-Hellman group~14 \cite{rfc2409,rfc4301}, ESP in
transport mode because GRE already provides the outer header. A tunnel \emph{interface} lets routing
protocols, QoS policies and firewall zones attach to it, and a broken tunnel is visible as a down
interface. GRE was chosen over a pure IPsec virtual tunnel interface because one GRE tunnel carries IPv4 and
IPv6 together; an IKEv1 \texttt{ipsec ipv4} tunnel carries IPv4 only and would leave the IPv6 BGP sessions
down. MTU and TCP MSS are lowered to leave room for the GRE and ESP headers.

\cfgsnippet[lst:ipsec]{ipsec.cfg}{IKE policy, IPsec profile and GRE tunnel interface (hub)}

\textbf{iBGP with route reflectors.} All six edge routers form AS~\FactAsTvi{}. Normal iBGP requires a full
mesh, because a router does not pass routes learned from one iBGP peer to another. The two HQ edge routers
are \emph{route reflectors} \cite{rfc4456} and the branch routers are their clients, so each branch router
needs a single session over its tunnel (\cref{lst:bgp-hub}). A reflector does not change the next hop of a
reflected route; a Da Nang router could not reach the tunnel address of a Hanoi router, so the hubs set
\texttt{next-hop-self all} -- \texttt{all} extends it to reflected routes. Both address families run over
the same tunnels: IPv4 sessions between the tunnel addresses and parallel IPv6 sessions using
multiprotocol BGP \cite{rfc4760}.

\cfgsnippet[lst:bgp-hub]{bgp_hub.cfg}{BGP on the HQ hub: eBGP policy and route-reflector clients}

\textbf{Path selection.} Each branch router tags the aggregate it sends to its hub with a local preference:
200 on the ISP-A router, 150 on the ISP-B router (\cref{lst:bgp-spoke}). Local preference is carried
through the whole AS, so every HQ router prefers the ISP-A tunnel to a branch and keeps the ISP-B tunnel as
backup. The traffic of a branch VLAN leaves through the tunnel of its active gateway, so paths may be
asymmetric; this is why the firewall rules between the LAN and the tunnels are stateless
(\cref{sec:security}).

\cfgsnippet[lst:bgp-spoke]{bgp_spoke.cfg}{BGP on a branch: local preference towards the hub}

\textbf{Failure detection.} Every iBGP session runs with 3\,s keepalive and 9\,s hold time. A tunnel has no
physical link to its peer, so the far end of a failed provider connection learns about it only when the
keepalives stop: 9\,s is therefore the dominant failover time on the WAN. The fast timers are also needed
on the session between the two HQ edge routers: their loopbacks lie inside the site's \texttt{Null0}
summary, so after one of them dies its loopback still ``resolves'' to the summary and only the hold timer
removes its routes (\cref{sec:challenges}). BFD \cite{rfc5880} would reduce detection below one second and is
the first item of the roadmap.

\textbf{Local Internet breakout.} Each edge router translates corporate sources to its WAN address (NAT
overload \cite{rfc3022}) for traffic to the Internet only: the NAT access list excludes
\texttt{10.0.0.0/8} destinations (\cref{lst:nat}). Inter-site traffic thus fails closed -- if the overlay
is down it is dropped, never leaked to the Internet in clear text -- and a Hanoi user reaches a SaaS service
through the Hanoi provider without crossing HQ, as the brief requires.

\subsection{Data-centre fabric: EVPN-VXLAN (design)}\label{sec:evpn}
This subsection is a design; the lab images do not support VXLAN or the EVPN address family.

The brief asks for a spine-and-leaf data centre with EVPN-VXLAN, and TVI's Hanoi AI workloads explain why:
model training moves large volumes between servers (east-west), which a spanning-tree network serves badly
because it blocks redundant links. In a spine-leaf fabric every leaf connects to every spine and all links
forward; any two servers are exactly two hops apart, and capacity grows by adding spines.

\begin{itemize}
  \item \textbf{Underlay:} routed point-to-point links between leaves and spines with eBGP (one AS per leaf)
        or OSPF, ECMP across all spines, loopbacks as VXLAN tunnel endpoints (VTEPs). Addresses come from the
        reserved block \texttt{10.1.128.0/18}.
  \item \textbf{Data plane:} VXLAN \cite{rfc7348} encapsulates Ethernet frames in UDP, so a segment can stretch
        across the routed fabric; its 24-bit identifier allows 16 million segments instead of 4{,}094 VLANs.
  \item \textbf{Control plane:} MP-BGP EVPN \cite{rfc7432,rfc8365} distributes MAC and IP reachability between
        VTEPs (spines act as route reflectors), replacing flood-and-learn and suppressing ARP flooding.
  \item \textbf{Gateways:} a distributed anycast gateway on every leaf, so a server's first hop is always its
        own leaf; L2VNIs for stretched segments, one L3VNI per tenant VRF (for example AI-lab, corporate
        applications) for routed traffic between them.
  \item \textbf{Integration:} a pair of border leaves connects to the campus core pair with routed links
        (VRF-lite) and advertises the data-centre prefixes into the HQ OSPF domain, which keeps the WAN design of
        this report unchanged.
\end{itemize}

\subsection{SD-WAN and wireless (design)}\label{sec:sdwan}
\textbf{SD-WAN.} Our overlay implements by hand what an SD-WAN controller automates: tunnels over two
providers, route exchange, path preference and local breakout. A controller-based solution would replace
the static tunnels with automatically built ones, measure loss, delay and jitter on every path, and steer
each application to the best path in real time (voice to the lowest-jitter path, backups to the cheapest),
with zero-touch provisioning of new branches. MPLS, which the brief lists as a second underlay, would
simply be a third transport in the same overlay. The addressing and the routing policy of this report
carry over unchanged.

\textbf{Wireless.} Lightweight access points would sit on dedicated access ports in the management VLAN and
build CAPWAP tunnels \cite{rfc5415} to a controller pair in the HQ data centre (stateful switchover). Users
would authenticate with 802.1X against RADIUS, which returns the VLAN of their role -- the dynamic VLAN
assignment the brief asks for -- while the guest SSID uses a captive portal and lands in VLAN~50. For the
autonomous guided vehicles in Da Nang, Wi-Fi 6/7 access points with fast roaming (802.11r/k/v) and a
dedicated, QoS-marked SSID mapped to the IoT automation VLAN would provide the required low latency.
